% ================================================================
% Anu Drive Methodology Template v1.0
% Auto-generated from series_registry.json
%
% Compile with: tectonic methodology.tex
% ================================================================
\documentclass[11pt,letterpaper]{article}

% --- Page layout ---
\usepackage[margin=1in]{geometry}

% --- Typography ---
\usepackage[T1]{fontenc}
\usepackage{lmodern}
\usepackage{microtype}

% --- Math (for formulas in construction steps) ---
\usepackage{amsmath,amssymb}

% --- Tables ---
\usepackage{booktabs}
\usepackage{longtable}
\usepackage{array}
\usepackage{tabularx}

% --- Colors ---
\usepackage[table]{xcolor}
\definecolor{headerblue}{HTML}{4472C4}
\definecolor{lightblue}{HTML}{D9E2F3}
\definecolor{lightgray}{HTML}{F2F2F2}

% --- Hyperlinks and PDF navigation ---
\usepackage[
    colorlinks=true,
    linkcolor=headerblue,
    urlcolor=headerblue,
    bookmarks=true,
    bookmarksnumbered=true,
    bookmarksopen=true,
    bookmarksopenlevel=2,
    pdfstartview=FitH
]{hyperref}
\usepackage{bookmark}

% --- Headers and footers ---
\usepackage{fancyhdr}
\pagestyle{fancy}
\fancyhf{}
\lhead{\small\textbf{<<PROJECT_NAME>>} --- Methodology}
\rhead{\small v<<VERSION>>}
\cfoot{\thepage}
\renewcommand{\headrulewidth}{0.4pt}

% --- Section formatting ---
\usepackage{titlesec}
\titleformat{\section}{\Large\bfseries}{\thesection}{1em}{}
\titleformat{\subsection}{\large\bfseries}{\thesubsection}{1em}{}
\titleformat{\subsubsection}{\normalsize\bfseries}{\thesubsubsection}{1em}{}

% --- Lists ---
\usepackage{enumitem}
\setlist[itemize]{nosep,left=0pt}
\setlist[enumerate]{nosep,left=0pt}
\setlist[description]{style=nextline,nosep}

% --- Paragraph spacing ---
\setlength{\parindent}{0pt}
\setlength{\parskip}{6pt}

% ================================================================
% DOCUMENT METADATA — Replace <<PLACEHOLDERS>> during generation
% ================================================================
\title{%
    \vspace{-1cm}
    {\Huge\bfseries <<PROJECT_TITLE>> \\[0.5cm]}
    {\Large Data Documentation \& Methodology \\[0.8cm]}
    {\large Version <<VERSION>> \\[0.3cm]}
    {\large <<DATE>>}
}
\author{%
    <<AUTHOR_NAME>> \\
    {\normalsize <<INSTITUTION>>} \\
    {\normalsize\texttt{<<EMAIL>>}}
}
\date{}

\begin{document}

% ================================================================
% TITLE PAGE
% ================================================================
\begin{titlepage}
\centering
\vspace*{3cm}

{\Huge\bfseries <<PROJECT_TITLE>> \\[0.8cm]}
{\Large Data Documentation \& Methodology \\[1.5cm]}
{\large Version <<VERSION>> \\[0.3cm]}
{\large <<DATE>> \\[3cm]}

{\Large <<AUTHOR_NAME>> \\[0.3cm]}
{\large <<INSTITUTION>> \\[0.5cm]}
{\normalsize \texttt{<<EMAIL>>} \\[4cm]}

\rule{0.6\textwidth}{0.4pt} \\[0.5cm]
{\normalsize Based on the empirical work in: \\[0.3cm]}
{\normalsize\itshape <<ORIGINAL_TITLE>> \\[0.2cm]}
{\normalsize by <<ORIGINAL_AUTHOR>> (<<ORIGINAL_YEAR>>) \\[0.2cm]}
{\normalsize <<PUBLISHER>>}

\vfill
\end{titlepage}

% ================================================================
% ABSTRACT
% ================================================================
\section*{Abstract}
\addcontentsline{toc}{section}{Abstract}

<<ABSTRACT_TEXT>>

% Pointer to package contents
\bigskip
\textbf{This package contains:}
\begin{itemize}
    \item \texttt{<<MASTER_XLSX>>} --- All <<N_SERIES>> final series in one spreadsheet
    \item \texttt{<<MASTER_CSV>>} --- Machine-readable CSV (for R, Python, Stata)
    \item \texttt{<<CODEBOOK_CSV>>} --- Formal data dictionary (one row per series)
    \item \texttt{Series/} --- Individual series workbooks with full construction detail
    \item \texttt{CITATION.txt} --- How to cite this data (formatted + BibTeX)
    \item This document --- Complete methodology for every series
\end{itemize}

\newpage

% ================================================================
% TABLE OF CONTENTS
% ================================================================
\tableofcontents
\newpage

% ================================================================
% CHAPTER 1: DATA OVERVIEW
% ================================================================
\section{Data Overview}
\label{sec:overview}

\subsection{Project Description}

<<PROJECT_DESCRIPTION>>

\subsection{Source Work}

This data package replicates and extends the empirical content of:

\begin{quote}
<<ORIGINAL_AUTHOR>>, \textit{<<ORIGINAL_TITLE>>} (<<ORIGINAL_YEAR>>). <<PUBLISHER>>.
\end{quote}

<<SOURCE_WORK_DESCRIPTION>>

\subsection{Coverage Summary}

\begin{tabular}{@{} l l @{}}
\toprule
\textbf{Dimension} & \textbf{Value} \\
\midrule
Time Coverage        & <<YEAR_START>> -- <<YEAR_END>> \\
Geographic Scope     & <<GEOGRAPHIC_SCOPE>> \\
Total Series         & <<N_SERIES>> \\
Extended to Present  & <<N_EXTENDED>> of <<N_SERIES>> \\
Unique Data Sources  & <<N_SOURCES>> \\
API Extensions       & <<API_LIST>> \\
\bottomrule
\end{tabular}

\subsection{Series Summary}
\label{sec:summary-table}

The following table lists all series in this package. Click any series ID to jump to its detailed documentation.

\begin{longtable}{@{} l p{4.5cm} c l l c l @{}}
\toprule
\textbf{ID} & \textbf{Name} & \textbf{Ch.} & \textbf{Period} & \textbf{Units} & \textbf{Ext.} & \textbf{Source} \\
\midrule
\endhead
<<SERIES_SUMMARY_ROWS>>
\bottomrule
\end{longtable}

\newpage

% ================================================================
% CHAPTER 2: SERIES DOCUMENTATION
% ================================================================
\section{Series Documentation}
\label{sec:series}

% --- BEGIN AUTO-GENERATED SERIES SECTIONS ---
% Each series gets one \subsection with standardized content.
% The generator produces one block per series from series_registry.json.

<<SERIES_SECTIONS>>

% --- END AUTO-GENERATED SERIES SECTIONS ---

\newpage

% ================================================================
% APPENDIX A: DATA SOURCES
% ================================================================
\appendix
\section{Data Sources}
\label{sec:sources}

Complete bibliography of all data sources used in this package, sorted alphabetically by institution.

<<DATA_SOURCES_BIBLIOGRAPHY>>

\newpage

% ================================================================
% APPENDIX B: FILE MANIFEST
% ================================================================
\section{File Manifest}
\label{sec:manifest}

\begin{longtable}{@{} p{6cm} p{8cm} @{}}
\toprule
\textbf{File} & \textbf{Description} \\
\midrule
\endhead
<<FILE_MANIFEST_ROWS>>
\bottomrule
\end{longtable}

\newpage

% ================================================================
% APPENDIX C: USING THIS DATA
% ================================================================
\section{Using This Data}
\label{sec:usage}

\subsection{Excel / Google Sheets}

Open \texttt{<<MASTER_XLSX>>} directly. The workbook is formatted for immediate use:

\begin{itemize}
    \item \textbf{Row 1}: Human-readable series names
    \item \textbf{Row 2}: Series identifiers (S001, S002, \ldots)
    \item \textbf{Row 3}: Units of measurement
    \item \textbf{Row 4+}: Data values (Year in Column A)
\end{itemize}

Headers are frozen so they remain visible while scrolling.

\subsection{R}

\begin{verbatim}
library(readr)

# Load all series (skip name and unit rows)
data <- read_csv("<<MASTER_CSV>>", skip = 2)

# Column names are series IDs: S001, S002, ...
# First column is Year
head(data)
\end{verbatim}

\subsection{Python}

\begin{verbatim}
import pandas as pd

# Load with series IDs as column names
data = pd.read_csv("<<MASTER_CSV>>", header=1, skiprows=[2])

# data.columns: ['Year', 'S001', 'S002', ...]
print(data.head())
\end{verbatim}

\subsection{Stata}

\begin{verbatim}
import delimited "<<MASTER_CSV>>", ///
    rowrange(4:) varnames(2) clear

describe
\end{verbatim}

\subsection{Understanding the Extenbooks}

Each file in the \texttt{Series/} folder is an Excel workbook with four sheets:

\begin{description}
    \item[Data] All raw data inputs (subsources) laid out as columns alongside every transformation and the final constructed series. This is full transparency: you see every number that went into the result.
    \item[Provenance] Where each data point came from: the source institution, the specific table or API, the time period covered, and quality classification.
    \item[Research] Quotes and methodology notes from the original author's text that guided the construction decisions.
    \item[Construction] A step-by-step log of operations: load, rebase, splice, extend. Each step shows its inputs and outputs.
\end{description}

\subsection{Codebook (Variable Reference)}

The file \texttt{<<CODEBOOK_CSV>>} is a formal data dictionary in CSV format with one row per series. It documents every variable in the master data file: series ID, name, definition, units, frequency, temporal coverage, observation count, data sources, and extension status. Import it alongside the master data file to programmatically look up variable metadata:

\begin{verbatim}
# Python
codebook = pd.read_csv("<<CODEBOOK_CSV>>")
codebook.set_index("series_id", inplace=True)
print(codebook.loc["S013", "definition"])
\end{verbatim}

\subsection{Citing This Data}

See \texttt{CITATION.txt} in the package root for a formatted citation (ready to paste into a manuscript) and a BibTeX entry (for \LaTeX\ users). Please also cite the original work.

\newpage

% ================================================================
% APPENDIX D: VERSION HISTORY
% ================================================================
\section{Version History}
\label{sec:versions}

\begin{tabular}{@{} l l p{10cm} @{}}
\toprule
\textbf{Version} & \textbf{Date} & \textbf{Changes} \\
\midrule
<<VERSION_HISTORY_ROWS>>
\bottomrule
\end{tabular}

\end{document}
